\documentclass[10pt]{article}

\usepackage[margin=0.75in]{geometry}
\usepackage{amsmath,amsthm,amssymb,bm}
\usepackage{xcolor}
\usepackage{cancel}
\usepackage{graphicx}
\usepackage{changepage}
\usepackage{circuitikz}
\usepackage{pgfplots}
\usepackage{physics}
\usepackage{tikz}
\usepackage{tikz-3dplot}
\usepackage{hyperref}
\usepackage{siunitx}
\usepackage{fontspec}
\usepackage{relsize}
\usepackage{subfig}
\usepackage{todonotes}
\usepackage{multicol, multirow, booktabs}
\usepackage[breakable]{tcolorbox}
\usepackage[inline]{enumitem}

\theoremstyle{definition}
\newtheorem{problem}{Problem}
\newtheorem{soln}{Solution}

\pgfplotsset{compat=newest}
\usetikzlibrary{lindenmayersystems}
\usetikzlibrary{arrows}
\usetikzlibrary{calc}
\usetikzlibrary{positioning, fit}
\usetikzlibrary{3d, perspective}
\usetikzlibrary{math}
\usetikzlibrary{fadings}
\usetikzlibrary{angles,quotes} 

\definecolor{incolor}{HTML}{303F9F}
\definecolor{outcolor}{HTML}{D84315}
\definecolor{cellborder}{HTML}{CFCFCF}
\definecolor{cellbackground}{HTML}{F7F7F7}
\newcommand{\ui}{\hat{i}}
\newcommand{\uj}{\hat{j}}
\newcommand{\uk}{\hat{k}}
\newcommand{\ux}{\hat{x}}
\newcommand{\uy}{\hat{y}}
\newcommand{\uz}{\hat{z}}
\newcommand{\uv}[1]{\hat{\bv{{#1}}}}
\newcommand{\pr}[1]{{#1}^\prime}
\pgfdeclarelayer{background}  
\pgfsetlayers{background,main}
\AtBeginDocument{\RenewCommandCopy\qty\SI}
\newcommand{\justif}[2]{&{#1}&\text{#2}}

\makeatletter
\newcommand{\boxspacing}{\kern\kvtcb@left@rule\kern\kvtcb@boxsep}
\makeatother
\newcommand{\prompt}[4]{
    \ttfamily\llap{{\color{#2}[#3]:\hspace{3pt}#4}}\vspace{-\baselineskip}
}

\newcommand{\thevenin}[2]{
  \begin{center}
    \begin{circuitikz} \draw
      (0,0) -- (2,0) to[battery1, l_=$V_{Th}\eq#1$] (2,2) 
      to[resistor, l_=$R_{Th}\eq#2$] (0,2)
      ;
      \draw [o-] (-.07,2.079);
      \draw [o-] (-.07,0.079);
    \end{circuitikz}
  \end{center}
}

\newcommand{\norton}[2]{
  \begin{center}
    \begin{circuitikz} \draw
      (0,0) -- (3,0) to[american current source, l_=$I_{N}\eq#1$] (3,2) -- (0,2) (2,0)
      to[resistor, l=$R_{N}\eq#2$] (2,2)
      ;
      \draw [o-] (-.07,2.079);
      \draw [o-] (-.07,0.079);
    \end{circuitikz}
  \end{center}
}

\DeclareMathOperator{\Div}{div}
\DeclareMathOperator{\Curl}{curl}

\DeclareSIUnit\mile{m}

\newlength\stdbls % standard baselineskip
\AtBeginDocument{%
  \setlength\stdbls{\baselineskip}%
  \gdef\stdblsn{\strip@pt\stdbls}%
} 

\newfontface{\Kaufmann}{Kaufmann}
\DeclareTextFontCommand{\kf}{\Kaufmann}
\newcommand{\scriptr}{\fontsize{12pt}{\stdbls}\kf{r}\fontsize{10pt}{\stdbls}}

\newfontface{\KaufmannB}{Kaufmann Bold}
\DeclareTextFontCommand{\kfb}{\KaufmannB}
\newcommand{\bscriptr}{\fontsize{12pt}{\stdbls}\kfb{r}\fontsize{10pt}{\stdbls}}

\newcommand{\highlight}[1]{\colorbox{yellow}{$\displaystyle #1$}}

\newcommand{\ti}[1]{\widetilde{#1}}

\newcommand{\bv}[1]{\bm{\mathrm{#1}}}

\title{Physics 3130H: Worksheet}
\author{Jeremy Favro (0805980) \\ Trent University, Peterborough, ON, Canada}
\date{\today}

\begin{document}
\maketitle

% PROBLEM 1
\begin{problem}
\begin{enumerate}[label=(\alph*)]
  \item Function representing curl is a \rule{1cm}{0.15mm}
  \item Function representing gradient is a \rule{1cm}{0.15mm}
  \item gravitational field is a \rule{1cm}{0.15mm}
  \item gravitational potential is a \rule{1cm}{0.15mm}
  \item forces are \rule{1cm}{0.15mm}
  \item force component is a \rule{1cm}{0.15mm}
  \item Speed is a \rule{1cm}{0.15mm}
  \item velocity is a \rule{1cm}{0.15mm}
  \item angular velocity is a \rule{1cm}{0.15mm}
  \item angular acceleration is a \rule{1cm}{0.15mm}
  \item potential energy is a \rule{1cm}{0.15mm}
  \item kinetic energy is a \rule{1cm}{0.15mm}
\end{enumerate}
\end{problem}
\begin{soln}~
  \begin{enumerate}[label=(\alph*)]
    \item Function representing curl is a \textbf{vector}.
    \item Function representing gradient is a \textbf{vector}.
    \item gravitational field is a \textbf{vector}.
    \item gravitational potential is a \textbf{scalar}.
    \item forces are \textbf{vectors}.
    \item force component is a \textbf{scalar}.
    \item speed is a \textbf{scalar}.
    \item velocity is a \textbf{vector}.
    \item angular velocity is a \textbf{vector}.
    \item angular acceleration is a \textbf{vector}.
    \item potential energy is a \textbf{scalar}.
    \item kinetic energy is a \textbf{scalar}.
  \end{enumerate}
\end{soln}
\newpage

% PROBLEM 2
\begin{problem}
An object moves in the $xy$-plane in such a way that its position vector $\bv{r}$ is given as a function
of time $t$ by $\bv{r}=a\cos \omega t \ui + b\sin\omega t\uj$ where $a$, $b$, and $\omega$ are constants.
\begin{enumerate}[label=(\alph*)]
  \item How far is the object from the origin at any time t?
  \item Find the object's velocity and acceleration as a function of time.
  \item Show that the object moves on the elliptical path.
\end{enumerate}
\end{problem}
\begin{soln}~
  \begin{enumerate}[label=(\alph*)]
    \item If we take the origin to be $\bv{O}=(0,0,0)$ then its distance from the origin is
          $r=\sqrt{a^2\cos^2 \omega t + b^2\sin^2\omega t}$.
    \item $\bv{v}=d\bv{r}/dt=-a\omega\sin \omega t \ui + b\omega\cos\omega t\uj$.
          $\bv{a}=d\bv{v}/dt=-a\omega^2\cos \omega t \ui - b\omega^2\sin\omega t\uj$.
    \item It is clear enough to see the elliptic motion which arises from these equations. The object's position ($r$, distance from the origin)
          is a strictly oscillating function, i.e. the object is confined to orbit around the origin.
  \end{enumerate}
\end{soln}

% PROBLEM 3
\begin{problem}
Calculate the divergence of each of the following functions:
\begin{enumerate}[label=(\alph*)]
  \item $x^2\ui+y^2\uj+z^2\uk$.
  \item $zy\ui+xz\uj+xy\uk$.
  \item $e^{-x}\ui+e^{-y}\uj+e^{-z}\uk$.
\end{enumerate}
\end{problem}
\begin{soln}~
  \begin{enumerate}[label=(\alph*)]
    \item $\bv{\nabla}\cdot(x^2\ui+y^2\uj+z^2\uk)=2x+2y+2z$.
    \item $\bv{\nabla}\cdot(zy\ui+xz\uj+xy\uk)=0$.
    \item $\bv{\nabla}\cdot(e^{-x}\ui+e^{-y}\uj+e^{-z}\uk)=-e^{-x}-e^{-y}-e^{-z}$.
  \end{enumerate}
\end{soln}

% PROBLEM 4
\begin{problem}
Calculate the curl of each of the following functions:
\begin{enumerate}[label=(\alph*)]
  \item $xy\ui+yz\uj-y^2\uk$.
  \item $xe^{-x}\ui+e^{-y}\uj+e^{-z}\uk$.
  \item $x\ui+y\uj+(x^2+y^2)\uk$.
\end{enumerate}
\end{problem}
\begin{soln}~
  \begin{enumerate}[label=(\alph*)]
    \item $\bv{\nabla}\times (xy\ui+yz\uj-y^2\uk)=\left[-2y-y\right]\ui+0\uj+x\uk$.
    \item $\bv{\nabla}\times (xe^{-x}\ui+e^{-y}\uj+e^{-z}\uk)=\bv{0}$.
    \item $\bv{\nabla}\times (x\ui+y\uj+(x^2+y^2)\uk)=2y\ui+2x\uj+0\uk$.
  \end{enumerate}
\end{soln}

% PROBLEM 5
\begin{problem}
For any central force show that curl of $\bv{F}$ is zero. (use spherical polar coordinates: Refer to
the lecture slides in gravitation chapter and central forces chapter)
\end{problem}
\begin{soln}
  Central forces are forces of the form $\bv{F}(\bv{r})=F(r)\uv{r}$. Curl in spherical coordinates is given by
  \begin{align*}
    \bv{\nabla}\times \bv{v} & =\frac{1}{r\sin\theta}\left[\partial_\theta(\sin\theta v_\phi) -\partial_\phi v_\theta\right]\uv{r} \\
                             & +\frac{1}{r}\left[\frac{1}{\sin\theta}\partial_\phi v_r-\partial_r(rv_\phi)\right]\uv{\theta}       \\
                             & +\frac{1}{r}\left[\partial_r(rv_\theta)-\partial_\theta v_r\right]\uv{\phi}.
  \end{align*}
  Clearly since the curl does not contain a term like $\partial_rv_r$ it wil be zero. Note that under this definition a central force is one
  which is spherically symmetric. A nonsymmetric central force which looks like $\bv{F}(\bv{r})=F(r,\theta,\phi)\uv{r}$ may not have
  vanishing curl.
\end{soln}

% PROBLEM 6
\begin{problem}
Objective is to find the
gravitational field. For the following configurations write down the corresponding
integrals:
\begin{enumerate}[label=(\alph*)]
  \item A rod of length $L$ (at a point $R$ along the length ($R>L/2$) from the center, at a point $R$
        perpendicular to the length from the center of the rod)
  \item A spherical shell (at a point inside the shell, in between the inner and outer radius, at a
        point outside the sphere)
  \item A hollow sphere (at a point inside the shell, and outside the shell)
  \item A solid sphere (at a point inside the shell, and outside the shell)
\end{enumerate}
\end{problem}
\begin{soln}
  Note that for the most part here I am just looking at setting up the integrals. I'm not going to make simplifications.
  \begin{enumerate}[label=(\alph*)]
    \item The general form for gravitational field is (using the notational conventions from Griffiths' electrodynamics)
          $$\bv{g}=-G\int\frac{\rho(\bv{r}')\uv{\bscriptr}}{\scriptr^2}d^3r'.$$
          Here we ought to take advantage of cylindrical symmetry and integrate over $s,\phi,z$. Additionally,
          we take this rod, given that we are not told a radius, to have a linear mass density $\lambda(\bv{r}')$.
          If we are at a point $R$ along the length of the rod then $\bv{r}=0\uv{s}+0\uv{\phi}+R\uz$ and
          $\bv{r}'=z'\uz$ so $\bv{\bscriptr}=R-z'\uz$ and the integral becomes
          $$\bv{g}=-G\int_0^L \frac{\lambda(z')\left[R-z'\right]}{\left|R-z'\right|^{3}}dz'\uz.$$
          When we are at a distance $R$ from the center of the rod along $\uv{s}$ we have that
          $\bv{r}=R\uv{s}$ and $\bv{r}'$ is unchanged. This gives $\bv{\bscriptr}=R\uv{s}-z'\uv{z}$ and so
          $$\bv{g}=-G\int_0^L \frac{\lambda(z')\left[R\uv{s}-z'\uv{z}\right]}{\left[R^2+(z')^2\right]^{3/2}}dz'.$$
    \item Inside the shell (of volume mass density $\bv{\rho}(\bv{r}')$), $\bv{r}=x\ux+y\uy+z\uz$,
          $\bv{r}'=x'\ux+y'\uy+z'\uz=r'\sin\theta'\cos\phi'\ux+r'\sin\theta'\sin\phi'\uy+r'\cos\theta'\uz$ and
          so
          $$\bv{\bscriptr}=(x-r'\sin\theta'\cos\phi')\ux+(y-r'\sin\theta'\sin\phi')\uy+(z-r'\cos\theta')\uz.$$
          The integral then becomes
          $$\bv{g}=-G\int_0^{2\pi}\int_0^\pi\int_0^R \frac{\rho(\bv{r}')\left[(x-r'\sin\theta'\cos\phi')\ux+(y-r'\sin\theta'\sin\phi')\uy+(z-r'\cos\theta')\uz\right]}{\left[(x-r'\sin\theta'\cos\phi')^2+(y-r'\sin\theta'\sin\phi')^2+(z-r'\cos\theta')^2\right]^{3/2}}(r')^2\sin\theta'dr'd\theta'd\phi'$$
          where I'm assuming the mass density will account for the hollow portion. If this is not the case the bounds on the $r'$ integral could be adjusted to integrate
          only over the real mass. Between the inner and outer shell the integral remains the same, the values of $x,y,z$ would just change.
          The same is the case for outside the sphere.
    \item I'll skimp on the details because I think I get the point. The integral is
          $$\bv{g}=-G\int_0^{2\pi}\int_0^\pi\frac{\sigma(\bv{r}')\left[(x-R\sin\theta'\cos\phi')\ux+(y-R\sin\theta'\sin\phi')\uy+(z-R\cos\theta')\uz\right]}{\left[(x-R\sin\theta'\cos\phi')^2+(y-R\sin\theta'\sin\phi')^2+(z-R\cos\theta')^2\right]^{3/2}}(R)^2\sin\theta'd\theta'd\phi'$$
          where I've assumed the shell has surface mass density $\sigma$ and radius $R$. Again this integral is the same inside and out, the values of $x,y,z$ will just change.
    \item This is the same integral as in part (b), just with $\rho(\bv{r}')$ nonzero for the entire sphere. In the case we would not have to change the bounds
          on the $r'$ integral as we are integrating over the entire spherical volume.
  \end{enumerate}
\end{soln}

% PROBLEM 7
\begin{problem}
Calculate the Coriolis force for an object falling vertically down in the southern
hemisphere. I am interested to see the components of the vector quantities involved, the
vector product, and the direction of the lateral deflection.
\end{problem}
\begin{soln}
  The Coriolis force is given by $-2m\bv{\omega}\times \bv{v}_r$ where
  $\bv{\omega}$ is the angular velocity of the rotating body which causes the force and
  $\bv{v}_r$ is the velocity of the falling object as measured in the rotating frame, which here is the surface of the Earth.
  Recall that in class we defined $\bv{\omega}=-\omega\cos\lambda\ux+\omega\sin\lambda\uz$ with $\ux$ pointing south, $\uy$ pointing east, and $\uz$ normal to the Earth's surface.
  $\omega$ is the angular \emph{speed} and $\lambda$ is the latitude, running from $+\pi/2$ at the north pole to $0$ at the equator and $-\pi/2$ at the south pole.
  Since we are in the southern hemisphere, $\abs{\lambda}=-\lambda$ so $\bv{\omega}=-\omega\cos\ux-\omega\sin\abs{\lambda}\uz$. Assuming the object starts
  from rest and begins falling (e.g. it wasn't thrown upwards, or if it was we start our measurement from the throw's apex) then $\bv{v}_r=-gt\uz$
  if we take gravity to be along $\uz$ alone which is not quite correct but is sufficient. This then gives a cross product of the form
  $$-\bv{\omega}\times\bv{v}_r=-\begin{vmatrix}
      \ux                & \uy & \uz                      \\
      -\omega\cos\lambda & 0   & -\omega\sin\abs{\lambda} \\
      0                  & 0   & -gt                      \\
    \end{vmatrix}=0\ux+gt\omega\cos\lambda\uy+0\uz$$
  so the deflection is entirely along $\uy$ and points east.
\end{soln}
\end{document}